% =============================================================================
%  A single artificial neuron.
%
%  Adapted from NNTikZ, Copyright (c) 2024 Fraser Love, MIT Licence:
%  https://github.com/fraserlove/nntikz -- re-coloured with this template's
%  palette and reduced to a bare fragment.
%
%  Roles match the other network figures: inputs tolblue, unit tolmagenta,
%  output tolorange, wiring black.
% =============================================================================
\begin{tikzpicture}[
    >=LaTeX,
    node/.style={circle, minimum width=2.25em, draw, thick},
    innode/.style={node, draw=tolblue,     fill=tolblue!12},
    midnode/.style={node, draw=tolmagenta, fill=tolmagenta!12},
    outnode/.style={node, draw=tolorange,  fill=tolorange!12},
    arrow/.style={-latex, semithick, draw=black},
    wlabel/.style={midway, fill=white, inner sep=1.5pt}
]

    % Inputs
    \foreach \x in {0,1,2}
        \node[innode] (x\x) at (0,-\x*1.25) {$x_{\x}$};
    \node[innode] (xm) at (0,-4*1.25) {$x_{n}$};
    \path (x2) -- (xm) node[pos=0.35, scale=2] {\vdots};

    % Summation unit and output
    \node[midnode] (z) at (2.5,-2.5) {$z$};
    \node[outnode, right of=z, node distance=2.5cm] (yhat) {$\hat{y}$};

    % Weighted inputs
    \draw[arrow] (x0) -- (z) node[wlabel] {$w_0$};
    \foreach \x in {1,...,2}
        \draw[arrow] (x\x) -- (z) node[wlabel] {$w_{\x}$};
    \draw[arrow] (xm) -- (z) node[wlabel] {$w_{n}$};

    % Activation
    \draw[arrow] (z) -- (yhat) node[wlabel] {$\sigma(z)$};

    \node[above of=z, node distance=1.2cm] {$z = \bm{w}^\top \bm{x}$};
\end{tikzpicture}
